\documentclass[12pt, a4paper]{article}

% ======================== PACOTES ========================
\usepackage[utf8]{inputenc}
\usepackage[T1]{fontenc}
\usepackage[brazil]{babel}
\usepackage{geometry}
\usepackage{graphicx}
\usepackage{hyperref}
\usepackage{listings}
\usepackage{xcolor}
\usepackage{booktabs}
\usepackage{array}
\usepackage{enumitem}
\usepackage{tikz}
\usepackage{titlesec}
\usepackage{fancyhdr}
\usepackage{parskip}
\usepackage{adjustbox}
\usepackage[htt]{hyphenat}
\usepackage{microtype}

\geometry{margin=2.5cm}
\setlength{\emergencystretch}{3em}

% ======================== CORES ========================
\definecolor{codegreen}{rgb}{0.25,0.49,0.38}
\definecolor{codegray}{rgb}{0.5,0.5,0.5}
\definecolor{codepurple}{rgb}{0.58,0,0.82}
\definecolor{codeblue}{rgb}{0.13,0.29,0.53}
\definecolor{backcolour}{rgb}{0.96,0.96,0.96}
\definecolor{unicampblue}{RGB}{0, 83, 138}

% ======================== LISTINGS ========================
\lstdefinestyle{javastyle}{
    backgroundcolor=\color{backcolour},
    commentstyle=\color{codegreen},
    keywordstyle=\color{codeblue}\bfseries,
    numberstyle=\tiny\color{codegray},
    stringstyle=\color{codepurple},
    basicstyle=\ttfamily\footnotesize,
    breakatwhitespace=false,
    breaklines=true,
    captionpos=b,
    keepspaces=true,
    numbers=left,
    numbersep=8pt,
    showspaces=false,
    showstringspaces=false,
    showtabs=false,
    tabsize=4,
    frame=single,
    rulecolor=\color{codegray},
    language=Java
}
\lstset{style=javastyle}

% ======================== HYPERREF ========================
\hypersetup{
    colorlinks=true,
    linkcolor=unicampblue,
    urlcolor=unicampblue,
    citecolor=unicampblue
}

% ======================== CABEÇALHO ========================
\pagestyle{fancy}
\fancyhf{}
\fancyhead[L]{\small SI405 --- Padrões de Projeto}
\fancyhead[R]{\small Faculdade de Tecnologia --- Unicamp}
\fancyfoot[C]{\thepage}
\renewcommand{\headrulewidth}{0.4pt}

% ======================== TÍTULO ========================
\title{
    \vspace{-1cm}
    \textbf{\Large Relatório Técnico} \\[0.3cm]
    \textbf{\LARGE Refatoração com Padrões Criacionais}
}
\author{}
\date{}

\begin{document}

% ======================== CAPA ========================
\maketitle
\vspace{-1.5cm}

\begin{center}
    \begin{tabular}{ll}
        \textbf{Disciplina:}  & SI405 --- Padrões de Projeto \\
        \textbf{Instituição:} & Faculdade de Tecnologia --- Unicamp \\
        \textbf{Data:}        & 12 de maio de 2026 \\
    \end{tabular}
\end{center}

\vspace{0.5cm}
\hrule
\vspace{0.8cm}

\tableofcontents

\newpage

% ================================================================
% SEÇÃO 1
% ================================================================
\section{Análise dos Principais Problemas Arquiteturais Identificados}

Após análise criteriosa do código legado, foram identificados os seguintes problemas arquiteturais:

\subsection{Dispersão da Lógica de Criação de Objetos}

Os geradores de relatório (\texttt{PdfRelatorioGenerator}, \texttt{CsvRelatorioGenerator}, \texttt{JsonRelatorioGenerator}) eram instanciados diretamente dentro dos serviços por meio de blocos \texttt{switch/case} e \texttt{if/else-if}:

\begin{itemize}
    \item Em \texttt{RelatorioService.gerarRelatorio()} --- cadeia de \texttt{if/else-if} com \texttt{new PdfRelatorioGenerator()}, \texttt{new CsvRelatorioGenerator()}, \texttt{new JsonRelatorioGenerator()}
    \item Em \texttt{ExportacaoService.exportar()} --- \texttt{switch/case} com a mesma lógica duplicada
\end{itemize}

Isso significava que \textbf{a lógica de criação estava replicada em dois locais distintos}, violando o princípio DRY (\textit{Don't Repeat Yourself}) e o Princípio da Responsabilidade Única (SRP).

\subsection{Alto Acoplamento entre Serviços e Implementações Concretas}

Tanto \texttt{RelatorioService} quanto \texttt{ExportacaoService} dependiam diretamente das classes concretas dos geradores. Não havia nenhuma abstração (interface ou classe abstrata) unificando os geradores, o que:

\begin{itemize}
    \item Impedia a substituição de implementações sem modificar os serviços
    \item Dificultava a adição de novos formatos (cada novo formato exigia alteração em múltiplos arquivos)
    \item Impossibilitava injeção de dependência e testes com mocks
\end{itemize}

\subsection{Instâncias Inconsistentes de Configuração}

O problema mais sutil e de maior impacto identificado: \textbf{cada serviço criava sua própria instância de \texttt{ConfiguracaoSistema} com valores diferentes}:

\begin{table}[h!]
\centering
\begin{adjustbox}{max width=\textwidth}
\begin{tabular}{lllll}
\toprule
\textbf{Serviço} & \textbf{Empresa} & \textbf{Ambiente} & \textbf{Diretório} & \textbf{Debug} \\
\midrule
\texttt{RelatorioService}    & ``Empresa XPTO''       & ``DEV''  & \texttt{/tmp/relatorios}       & \texttt{false} \\
\texttt{ExportacaoService}   & ``Empresa XPTO Ltda.'' & ``PROD'' & \texttt{/var/exports/relatorios} & \texttt{false} \\
\texttt{ConfiguracaoService} & ``Empresa XPTO Ltda.'' & ``DEV''  & \texttt{/tmp/relatorios}       & \texttt{true}  \\
\bottomrule
\end{tabular}
\end{adjustbox}
\caption{Configurações divergentes entre os serviços no código legado}
\end{table}

Isso violava a premissa de que configurações globais devem ser \textbf{únicas e consistentes} em toda a aplicação. Uma alteração de ambiente em um serviço não se refletia nos demais.

\subsection{Falta de Extensibilidade}

A inexistência de uma interface comum para os geradores e a dispersão da criação por \texttt{switch/case} tornavam a adição de um novo formato de relatório um processo trabalhoso e propenso a erros, exigindo alterações em pelo menos 3 arquivos.

% ================================================================
% SEÇÃO 2
% ================================================================
\section{Justificativa das Decisões Arquiteturais}

\subsection{Por que Factory Method?}

O padrão \textbf{Factory Method} foi escolhido porque os problemas identificados (duplicação de lógica de criação, alto acoplamento com classes concretas, dificuldade de extensão) são exatamente os cenários clássicos que esse padrão resolve:

\begin{itemize}
    \item \textbf{Centraliza a decisão de criação} em um único local (\texttt{RelatorioGeneratorFactory})
    \item \textbf{Desacopla os serviços} das implementações concretas por meio da interface \texttt{RelatorioGenerator}
    \item \textbf{Facilita extensão}: adicionar um novo formato requer apenas criar a classe geradora e adicionar uma linha na fábrica
\end{itemize}

\subsection{Por que Singleton?}

O padrão \textbf{Singleton} foi escolhido para \texttt{ConfiguracaoSistema} porque:

\begin{itemize}
    \item A configuração do sistema representa um \textbf{recurso compartilhado} que deve ser globalmente consistente
    \item Múltiplas instâncias com valores divergentes é um defeito real (já presente no código legado)
    \item O Singleton garante \textbf{ponto único de acesso} e \textbf{unicidade da instância}
\end{itemize}

\subsection{Decisão de Manter o Construtor Público}

O construtor público de \texttt{ConfiguracaoSistema} foi mantido por dois motivos:

\begin{enumerate}
    \item \textbf{Compatibilidade retroativa}: os testes legados criam instâncias independentes para testar getters/setters
    \item \textbf{Flexibilidade}: permite cenários de teste onde configurações isoladas são necessárias
\end{enumerate}

A unicidade é garantida pelo método \texttt{getInstancia()} para uso na aplicação, enquanto o construtor serve para testes e cenários excepcionais.

% ================================================================
% SEÇÃO 3
% ================================================================
\section{Explicação da Aplicação do Factory Method}

\subsection{Estrutura Implementada}

A Figura~\ref{fig:factory-diagram} apresenta o diagrama de classes da estrutura implementada com o padrão Factory Method.

\begin{figure}[h!]
\centering
\begin{adjustbox}{max width=\textwidth}
\begin{tikzpicture}[
    class/.style={rectangle, draw=black, fill=backcolour, minimum width=4.8cm, minimum height=1.2cm, text centered, font=\small\ttfamily},
    interface/.style={rectangle, draw=unicampblue, fill=blue!5, minimum width=4.8cm, minimum height=1.6cm, text centered, font=\small\ttfamily},
    factory/.style={rectangle, draw=black!70, fill=yellow!10, minimum width=5cm, minimum height=1.2cm, text centered, font=\small\ttfamily},
    arrow/.style={->, >=stealth, thick, dashed},
    uses/.style={->, >=stealth, thick}
]

% Interface
\node[interface] (iface) at (0, 4) {
    \begin{tabular}{c}
    \textit{<<interface>>} \\
    \textbf{RelatorioGenerator} \\
    \scriptsize +gerar(Relatorio): String
    \end{tabular}
};

% Factory
\node[factory] (factory) at (8, 4) {
    \begin{tabular}{c}
    \textbf{RelatorioGeneratorFactory} \\
    \scriptsize +criarGerador(Formato): Generator
    \end{tabular}
};

% Concrete classes
\node[class] (pdf) at (-4, 0.5) {\textbf{PdfRelatorioGenerator}};
\node[class] (csv) at (-4, -0.8) {\textbf{CsvRelatorioGenerator}};
\node[class] (json) at (0, -0.8) {\textbf{JsonRelatorioGenerator}};
\node[class] (xml) at (4, -0.8) {\textbf{XmlRelatorioGenerator}};
\node[class] (html) at (4, 0.5) {\textbf{HtmlRelatorioGenerator}};

% Arrows: implements
\draw[arrow] (pdf) -- (iface);
\draw[arrow] (csv) -- (iface);
\draw[arrow] (json) -- (iface);
\draw[arrow] (xml) -- (iface);
\draw[arrow] (html) -- (iface);

% Arrow: factory creates
\draw[uses] (factory) -- node[above, font=\scriptsize] {cria} (iface);

\end{tikzpicture}
\end{adjustbox}
\caption{Diagrama de classes --- Padrão Factory Method aplicado aos geradores de relatório}
\label{fig:factory-diagram}
\end{figure}

\subsection{Componentes Criados}

\begin{enumerate}
    \item \textbf{\texttt{RelatorioGenerator}} (interface) --- contrato comum com o método \texttt{gerar(Relatorio): String}
    \item \textbf{\texttt{RelatorioGeneratorFactory}} --- fábrica centralizada com o método \texttt{criarGerador(FormatoRelatorio)}
    \item \textbf{\texttt{XmlRelatorioGenerator}} --- novo gerador com escape de caracteres especiais XML
    \item \textbf{\texttt{HtmlRelatorioGenerator}} --- novo gerador com estrutura HTML5 completa
\end{enumerate}

\subsection{Refatoração nos Serviços}

\textbf{Antes} (\texttt{RelatorioService}):

\begin{lstlisting}
if (formato == FormatoRelatorio.PDF) {
    PdfRelatorioGenerator generator = new PdfRelatorioGenerator();
    return generator.gerar(relatorio);
} else if (formato == FormatoRelatorio.CSV) {
    CsvRelatorioGenerator generator = new CsvRelatorioGenerator();
    return generator.gerar(relatorio);
} else if (formato == FormatoRelatorio.JSON) {
    // ...
}
\end{lstlisting}

\textbf{Depois} (\texttt{RelatorioService}):

\begin{lstlisting}
RelatorioGenerator generator = generatorFactory.criarGerador(formato);
return generator.gerar(relatorio);
\end{lstlisting}

A mesma simplificação foi aplicada ao \texttt{ExportacaoService}, eliminando completamente o bloco \texttt{switch/case}.

% ================================================================
% SEÇÃO 4
% ================================================================
\section{Explicação da Aplicação do Singleton}

\subsection{Estrutura Implementada}

A Figura~\ref{fig:singleton-diagram} apresenta o diagrama de classes da estrutura implementada com o padrão Singleton.

\begin{figure}[h!]
\centering
\begin{adjustbox}{max width=\textwidth}
\begin{tikzpicture}[
    singleton/.style={rectangle, draw=unicampblue, fill=blue!5, minimum width=6.2cm, minimum height=3cm, text centered, font=\small\ttfamily},
    service/.style={rectangle, draw=black, fill=backcolour, minimum width=4cm, minimum height=1cm, text centered, font=\small\ttfamily},
    arrow/.style={->, >=stealth, thick}
]

% Singleton class
\node[singleton] (config) at (0, 3) {
    \begin{tabular}{c}
    \textbf{ConfiguracaoSistema} \\[4pt]
    \scriptsize --static instancia: ConfiguracaoSistema \\
    \scriptsize --nomeEmpresa: String \\
    \scriptsize --ambiente: String \\
    \scriptsize --diretorioExportacao: String \\
    \scriptsize --debugAtivo: boolean \\[4pt]
    \scriptsize +static getInstancia(): ConfiguracaoSistema \\
    \scriptsize +static resetInstancia(): void
    \end{tabular}
};

% Services
\node[service] (rs) at (-5, -1) {\textbf{RelatorioService}};
\node[service] (es) at (0, -1) {\textbf{ExportacaoService}};
\node[service] (cs) at (5, -1) {\textbf{ConfiguracaoService}};

% Arrows
\draw[arrow] (rs) -- node[left, font=\scriptsize, pos=0.5] {getInstancia()} (config);
\draw[arrow] (es) -- node[left, font=\scriptsize, pos=0.5] {getInstancia()} (config);
\draw[arrow] (cs) -- node[right, font=\scriptsize, pos=0.5] {getInstancia()} (config);

\end{tikzpicture}
\end{adjustbox}
\caption{Diagrama de classes --- Padrão Singleton aplicado à configuração do sistema}
\label{fig:singleton-diagram}
\end{figure}

\subsection{Implementação}

\begin{itemize}
    \item \textbf{Campo estático privado} \texttt{instancia} armazena a referência única
    \item \textbf{Inicialização preguiçosa} (\textit{lazy initialization}) no método \texttt{getInstancia()}
    \item \textbf{Método \texttt{resetInstancia()}} para permitir isolamento em testes automatizados
    \item \textbf{Construtor público preservado} para compatibilidade com testes legados
\end{itemize}

\subsection{Resultado}

\textbf{Antes:} três instâncias independentes com valores divergentes (nomes de empresa diferentes, ambientes diferentes, diretórios diferentes).

\textbf{Depois:} todos os serviços compartilham exatamente a mesma instância via \texttt{ConfiguracaoSistema.getInstancia()}, garantindo consistência global.

% ================================================================
% SEÇÃO 5
% ================================================================
\section{Impactos da Refatoração na Organização e Extensibilidade do Sistema}

\subsection{Resumo Quantitativo}

\begin{table}[h!]
\centering
\begin{adjustbox}{max width=\textwidth}
\begin{tabular}{lll}
\toprule
\textbf{Métrica} & \textbf{Antes} & \textbf{Depois} \\
\midrule
Formatos suportados             & 3 (PDF, CSV, JSON)   & 5 (+ XML, HTML)            \\
Pontos de criação de geradores  & 2 (duplicados)       & 1 (fábrica centralizada)   \\
Instâncias de configuração      & 3 (inconsistentes)   & 1 (Singleton)              \\
Testes automatizados            & 23                   & 60                         \\
Arquivos de teste               & 3                    & 6                          \\
\bottomrule
\end{tabular}
\end{adjustbox}
\caption{Comparativo quantitativo antes e depois da refatoração}
\end{table}

\subsection{Impacto na Extensibilidade}

\textbf{Adicionar um novo formato de relatório agora requer apenas:}

\begin{enumerate}
    \item Criar a classe geradora implementando \texttt{RelatorioGenerator}
    \item Adicionar o valor ao enum \texttt{FormatoRelatorio}
    \item Adicionar uma linha na fábrica \texttt{RelatorioGeneratorFactory}
\end{enumerate}

Nenhum serviço (\texttt{RelatorioService}, \texttt{ExportacaoService}) precisa ser modificado, pois eles dependem apenas da interface.

\subsection{Impacto na Manutenção}

\begin{itemize}
    \item \textbf{Redução de duplicação}: a lógica de criação de geradores existe em apenas um local
    \item \textbf{Consistência de configuração}: impossibilidade de ter ambientes divergentes entre serviços
    \item \textbf{Código mais legível}: os serviços delegam a criação e focam na lógica de negócio
    \item \textbf{Testabilidade}: a interface permite substituição por mocks em testes unitários
\end{itemize}

\subsection{Princípios SOLID Aplicados}

\begin{table}[h!]
\centering
\begin{adjustbox}{max width=\textwidth}
\begin{tabular}{>{\bfseries}l p{10cm}}
\toprule
\textbf{Princípio} & \textbf{Aplicação} \\
\midrule
SRP & A fábrica é responsável exclusivamente pela criação; os serviços focam na lógica de negócio \\
OCP & Novos formatos podem ser adicionados sem modificar os serviços existentes \\
DIP & Os serviços dependem da abstração \texttt{RelatorioGenerator}, não das classes concretas \\
\bottomrule
\end{tabular}
\end{adjustbox}
\caption{Princípios SOLID aplicados na refatoração}
\end{table}

\subsection{Testes Adicionados}

\begin{itemize}
    \item \textbf{\texttt{RelatorioGeneratorFactoryTest}} (22 testes) --- comportamento da fábrica, geradores XML e HTML, integração
    \item \textbf{\texttt{SingletonConfiguracaoTest}} (8 testes) --- unicidade, valores padrão, reset, compartilhamento entre serviços
    \item \textbf{\texttt{ExportacaoNovosFormatosTest}} (7 testes) --- exportação nos formatos XML e HTML
\end{itemize}

% ================================================================
% SEÇÃO 6
% ================================================================
\section{Arquivos Modificados e Criados}

\subsection{Arquivos Criados}

\begin{itemize}
    \item \texttt{generator/RelatorioGenerator.java} --- interface dos geradores
    \item \texttt{generator/RelatorioGeneratorFactory.java} --- fábrica centralizada
    \item \texttt{generator/XmlRelatorioGenerator.java} --- gerador XML
    \item \texttt{generator/HtmlRelatorioGenerator.java} --- gerador HTML
    \item \texttt{test/.../RelatorioGeneratorFactoryTest.java} --- testes da fábrica e novos formatos
    \item \texttt{test/.../SingletonConfiguracaoTest.java} --- testes do Singleton
    \item \texttt{test/.../ExportacaoNovosFormatosTest.java} --- testes de exportação novos formatos
\end{itemize}

\subsection{Arquivos Modificados}

\begin{itemize}
    \item \texttt{domain/ConfiguracaoSistema.java} --- adição do padrão Singleton
    \item \texttt{domain/FormatoRelatorio.java} --- adição de XML e HTML
    \item \texttt{generator/PdfRelatorioGenerator.java} --- implementa \texttt{RelatorioGenerator}
    \item \texttt{generator/CsvRelatorioGenerator.java} --- implementa \texttt{RelatorioGenerator}
    \item \texttt{generator/JsonRelatorioGenerator.java} --- implementa \texttt{RelatorioGenerator}
    \item \texttt{service/RelatorioService.java} --- uso da fábrica e Singleton
    \item \texttt{service/ExportacaoService.java} --- uso da fábrica e Singleton
    \item \texttt{service/ConfiguracaoService.java} --- uso do Singleton
    \item \texttt{Main.java} --- menu atualizado com novos formatos
\end{itemize}

\end{document}
